\section{Reduction to a compact box}\label{sec:reduction}

The class $\Ctwo$ is not compact: the period lattice can degenerate and the two cosets can approach each other. Both degenerations
cost energy, and this confines every possible violation of the main inequality to an explicit box.

\begin{lemma}[Separation]\label{lem:separation}
Let $Y=R\,\Z^3\cup(R\,\Z^3+Rf)$ be a normalized two-site configuration and $X=\ltwo Y$. If $|Rm|\le\frac58$ for some
$m\in\Z^3\setminus\{0\}$, or $|R(m+f)|\le\frac12$ for some $m\in\Z^3$, then $e(X)>\eF$.
\end{lemma}

\begin{proof}
By \eqref{eq:H2}, $e(X)>\eF$ is equivalent to $\Htwo(M,f)>4+2\ltwo\eF=0.36099256\ldots$. All terms of $\Htwo$ are nonnegative
and $\Ga(\frac12,\cdot)$ is decreasing. In the first case the same-coset terms $\pm m$ give
$\Htwo\ge2\,\Ga(\frac12,\frac{25\pi}{64})=0.37503436\ldots$, and in the second case one cross-coset term gives
$\Htwo\ge\Ga(\frac12,\frac\pi4)=0.42018281\ldots$.
\end{proof}

\begin{lemma}[Reduced period matrix]\label{lem:reduced}
Every lattice $L\subset\R^3$ of covolume one has, after an orthogonal change of coordinates, a basis matrix
\begin{equation}\label{eq:Rchart}
  R=\begin{pmatrix}r_1&r_1x&r_1y\\0&r_2&r_2z\\0&0&r_3\end{pmatrix},\qquad r_1,r_2>0,\quad r_3=\frac1{r_1r_2},\qquad
  x,y\in[0,\tfrac12],\quad z\in[-\tfrac12,\tfrac12],
\end{equation}
such that
\begin{align}
  |Rm|^2&\ge r_1^2 &&\text{for all }m\in\Z^3\setminus\{0\},\label{eq:R1}\\
  r_2^2(m_2+zm_3)^2+r_3^2m_3^2&\ge r_2^2 &&\text{for all }(m_2,m_3)\in\Z^2\setminus\{0\}.\label{eq:R2}
\end{align}
Consequently
\begin{equation}\label{eq:R3}
  r_2\ge\tfrac{\sqrt3}2\,r_1,\qquad r_3\ge\tfrac{\sqrt3}2\,r_2,\qquad r_1\le\tfrac2{\sqrt3},\qquad r_2\le\Bigl(\frac2{\sqrt3\,r_1}\Bigr)^{1/2}.
\end{equation}
\end{lemma}

\begin{proof}
This is the reduction of Hermite and of Korkine and Zolotarev \cite{KorkineZolotareff1873}. Let $b_1$ be a shortest nonzero vector of
$L$; it is primitive. The orthogonal projection $\pi(L)$ of $L$ onto $b_1^\perp$ is a lattice of rank two. Let $p$ be a shortest nonzero
vector of $\pi(L)$, complete it to a basis $(p,p')$ of $\pi(L)$, and choose $b_2,b_3\in L$ with $\pi(b_2)=p$, $\pi(b_3)=p'$. Then
$(b_1,b_2,b_3)$ is a basis of $L$. Subtracting integer multiples of $b_1$ from $b_2$, then of $b_2$ and of $b_1$ from $b_3$, we obtain
$|x|,|y|,|z|\le\frac12$ in the coordinates in which $b_1\parallel e_1$ and $b_2\in\operatorname{span}(e_1,e_2)$; this changes neither
$b_1$ nor $\pi(b_2)$. Replacing $b_2$ by $-b_2$ and reflecting the second coordinate maps $(x,y,z)$ to $(-x,y,-z)$, and replacing $b_3$
by $-b_3$ and reflecting the third coordinate maps it to $(x,-y,-z)$; so $x,y\ge0$ can be arranged. Now \eqref{eq:R1} says that $b_1$
is shortest, and \eqref{eq:R2} that $\pi(b_2)$ is shortest in $\pi(L)$. Taking $m=e_2$ in \eqref{eq:R1} gives
$r_2^2\ge(1-x^2)r_1^2\ge\frac34r_1^2$, and $(m_2,m_3)=(0,1)$ in \eqref{eq:R2} gives $r_3^2\ge\frac34r_2^2$. Then
$1=r_1r_2r_3\ge\frac{3\sqrt3}8r_1^3$ and $r_2^2\le\frac2{\sqrt3}r_2r_3=2/(\sqrt3\,r_1)$.
\end{proof}

We combine the five shape parameters of \eqref{eq:Rchart} and the fractional offset into one point $\kappa=(r_1,r_2,x,y,z,f)$ of
$\R^8$, and we write $R(\kappa)$, $Y(\kappa)=R(\kappa)\Z^3\cup(R(\kappa)\Z^3+R(\kappa)f)$ and $X(\kappa)=\ltwo Y(\kappa)$. Let
\begin{equation}\label{eq:Kbox}
  \Kbox:=\bigl[\tfrac58,\tfrac{37}{32}\bigr]\times\bigl[\tfrac12,\tfrac{11}8\bigr]\times\bigl[0,\tfrac12\bigr]^2
      \times\bigl[-\tfrac12,\tfrac12\bigr]\times\bigl[-\tfrac12,\tfrac12\bigr]^3\subset\R^8 ,
\end{equation}
and let $\Kbox_{\mathrm{red}}$ be the set of $\kappa\in\Kbox$ with $f\ne0$ for which \eqref{eq:R1} and \eqref{eq:R2} hold. On
$\Kbox_{\mathrm{red}}$ we regard $e$, $D$ and $\Def$ as functions of $\kappa$.

\begin{proposition}[Compact reduction]\label{prop:compact}
Let $X\in\Ctwo$ with $e(X)\le\eF$; in particular, let $X\in\Ctwo$ with $\Def(X)\le0$. Then there are $Q\in\Or(3)$, $v\in\R^3$ and
$\kappa\in\Kbox_{\mathrm{red}}$ with $X=QX(\kappa)+v$. Moreover the normalized configuration $Y(\kappa)$ satisfies
\begin{equation}\label{eq:sepstrict}
  |R(\kappa)m|>\tfrac58\ \ (m\in\Z^3\setminus\{0\}),\qquad |R(\kappa)(m+f)|>\tfrac12\ \ (m\in\Z^3).
\end{equation}
\end{proposition}

\begin{proof}
If $\Def(X)\le0$ then $e(X)\le\eB+\gap D(X)\le\eF$. Translate one point of $X$ to the origin and apply Lemma~\ref{lem:reduced} to the
normalized period lattice; the offset is transformed with the basis and then reduced modulo $\Z^3$ to $f\in[-\frac12,\frac12]^3$, which
changes neither $Y$ nor, by Lemma~\ref{lem:Hsym}, the energy. Lemma~\ref{lem:separation} gives \eqref{eq:sepstrict}; in particular
$r_1>\frac58$ and $f\ne0$. With \eqref{eq:R3},
\[
  r_1\le\tfrac2{\sqrt3}<\tfrac{37}{32},\qquad r_2\ge\tfrac{5\sqrt3}{16}>\tfrac12,\qquad r_2<\Bigl(\frac{16}{5\sqrt3}\Bigr)^{1/2}<\tfrac{11}8 .\qedhere
\]
\end{proof}

Thus Theorem~\ref{thm:main} is equivalent to the statement that $\Def\ge0$ on $\Kbox_{\mathrm{red}}$, with equality only at the
points that represent the body-centred and face-centred cubic lattices. Three remarks on the chart $\Kbox$ are in order.

\emph{Redundancy.} The chart is far from injective: a configuration has many reduced representatives (ties in \eqref{eq:R1},
\eqref{eq:R2}, and the different descriptions of a lattice as a two-site configuration), and $\Kbox$ also contains points that are not
reduced. Volumes in $\Kbox$ therefore have no invariant meaning, and we never use them. What matters is that every configuration with
$e\le\eF$ has \emph{at least one} representative in $\Kbox_{\mathrm{red}}$, so that it suffices to prove the inequality at every point
of $\Kbox_{\mathrm{red}}$; a region of $\Kbox$ may be discarded as soon as it is shown to be disjoint from $\Kbox_{\mathrm{red}}$, that is,
as soon as one of the inequalities \eqref{eq:R1}, \eqref{eq:R2} is strictly violated throughout it for a fixed integer vector.

\emph{Reflection.} The exchange of the two cosets is $f\mapsto-f$ at fixed shape, and it preserves $e$ and $D$ (Lemmas~\ref{lem:Hsym}
and~\ref{lem:Dbasic}). It is therefore enough to treat the half $\Kbox^+:=\Kbox\cap\{f_1\ge0\}$.

\emph{The equality points.} Each of the lattices $\ltwo\Bz$ and $\ltwo\Fz$ has seven descriptions as a two-site configuration, one for
each sublattice of index two, and each description has finitely many reduced representatives in $\Kbox$; all of them have
$f\in\frac12\Z^3$. By Theorem~\ref{thm:main} these will turn out to be the only points of $\Kbox_{\mathrm{red}}$ where $\Def=0$; near
them $\Def$ is small. The cover treats them by the anchor rules of Section~\ref{ssec:anchor} (Remark~\ref{rem:margins}), and
Section~\ref{sec:local} gives explicit neighbourhoods of them in which \eqref{eq:main} holds.
